% Proposed isolated A.87. All checkpoint28 sources remain unchanged.
\statement{A.87}{Proved}{The actual marked three-cycle and the explicit canonical cusp frame}
Retain the actual negative involution \(\alpha=(-1,f)\in\widehat{GT}_0\),
its canonical \(u\in F'\), and the notation of A.85--A.86.
Use the actual modular marking
\[
 \begin{gathered}
 B=C_2(s)\amalg C_3(r),\qquad \delta=sr,\qquad
 x=\delta^{-2},\quad y=sxs,\quad z=(xy)^{-1},\\
 \rho=\operatorname{Ad}(r)|_F,\qquad
 \iota=\operatorname{Ad}(s)|_F,\qquad
 v=u^{-1}s,\qquad \tau=v\alpha.
 \end{gathered}
 \tag{A.87.1}
\]
Here \(\rho\) sends x,y,z cyclically to y,z,x. Automorphisms and
their involutory symbols in the indicated actual semidirect products
are distinguished by context. Let \(d=d_\alpha\) be A.45's full unit,
put \(m=(d-1)/2\), and set
\[
 T_d=v\delta^d,\qquad \Xi_d=T_d^3.
 \tag{A.87.2}
\]
Then the full marked cusp obstruction of A.86 has the exact equivalent
form
\[
                 \Xi_d=1\quad\Longleftrightarrow\quad\Omega_\alpha=1.
 \tag{A.87.3}
\]
Unconditionally \(\Xi_d\in F'\), \(\tau(\Xi_d)=\Xi_d^{-1}\), and
there is a unique \(a_d\in F\) with
\[
 \Xi_d=a_d\tau(a_d)^{-1},\qquad
 (a_d)_{\rm ab}=(\ell_d,-\ell_d).
 \tag{A.87.4}
\]
No vanishing of \(\Xi_d\) or \(a_d\) is asserted.

Let A be the unique pure r-frame of the \emph{actual} positive element
\(\Sigma=\alpha c=(1,f)\):
\[
 \Phi_\Sigma(r)=ArA^{-1}=f^{-1}r,\qquad A\in F.
 \tag{A.87.5}
\]
Then \(A\in F'\), and the unique normalized cusp frame of A.86 is
\[
 \begin{gathered}
 f=\rho(A)A^{-1},\qquad
 \alpha(A)=\iota(A^{-1})f,\qquad
 \tau(A)=u^{-1}\rho(A)^{-1}u,\\
 h_0=\rho^2(A)\rho(A)^{-1}u,\qquad
 h_0\tau=r^{-1}(\delta^{-1}\alpha)r.
 \end{gathered}
 \tag{A.87.6}
\]
Writing \(h_0=q\tau(q)^{-1}\) with A.86's unique q, and
\(q_{\rm ab}=(t,-t)\), its unique Alexander solution can be transported
to the original first-cusp marking. There is a unique \(G\in F\) with
\[
 \Phi_\alpha(G)=\delta^{-1}G\delta^{-1},\qquad
 G_{\rm ab}=(-2t,a_{\rm reg}),\qquad a_{\rm reg}=1-4t.
 \tag{A.87.7}
\]
Its image in \(F/\langle\!\langle x\rangle\!\rangle\), with generator y,
is \(y^{a_{\rm reg}}\). If \(k\) is the actual A.45 dihedral exponent,
the remaining coefficient comparison is precisely
\[
             a_{\rm reg}=d_\alpha
             \quad\Longleftrightarrow\quad G_{{\rm ab},x}=k.
 \tag{A.87.8}
\]
Neither this equality nor the nonlinear equation (A.87.3) is proved here.

\begin{proof}
\emph{1. Existing modular maps and their ordered identities.}\par
The actual multiplier-one map of \(\Sigma=\alpha c\) in the attached
LM7\_PUSH D.50.23 and E.65.8 is
\[
 \Phi_\Sigma(s)=sf,\qquad \Phi_\Sigma(r)=f^{-1}r,
 \qquad\Phi_\Sigma(\delta)=\delta.
\]
The standard c fixes s and inverts \(\delta\). Composing these existing
maps gives
\[
 \Phi_\alpha(s)=sf,\qquad
 \Phi_\alpha(\delta)=\delta^{-1},\qquad
 \Phi_\alpha(r)=sf r^{-1}s.
 \tag{A.87.9}
\]
The order-three relation in (A.87.9) is already valid because this is
the modular extension of the given gentle element. It is not inferred
from a prospective positive conjugator. The canonical equality
\(u=f\alpha(u)\) and \(\alpha^2=1\) give \(\alpha(f)=f^{-1}\).
Consequently (A.87.9) squares to the identity on s and \(\delta\),
and hence on B. The two-cycle relation is \(fsf=s\), so
\[
 \Phi_\alpha(v)=\alpha(u)^{-1}sf=u^{-1}fsf=v.
\]
Also \(v^2=1\), because \(\iota(u)=u^{-1}\). Thus \(\tau=v\alpha\)
is an actual involution of B, whose restriction to F is exactly
the regular involution in A.86. Finally the actual half rotation is
\[
 e=(xu)v=xs=r^{-1}\delta r,\qquad e^2=xy=z^{-1}.
 \tag{A.87.10}
\]
This follows by reducing the words with \(s^2=r^3=1\).

\emph{2. The ordered hexagon and its exact scope.}\par
Put \(u_0=u(x,y)\), \(u_1=u(y,z)\), and \(u_2=u(z,x)\).
Since d is odd, \(\delta^d=\delta x^{-m}\), and direct multiplication gives
\[
 \begin{aligned}
 T_d&=u_0^{-1}r x^{-m}=u_0^{-1}y^{-m}r,\\
 \Xi_d&=u_0^{-1}y^{-m}u_1^{-1}z^{-m}u_2^{-1}x^{-m},\\
 \Xi_d^{-1}&=x^m u_2z^m u_1y^m u_0.
 \end{aligned}
 \tag{A.87.11}
\]
The power identities extend from integers to profinite m by continuity.
The S3-image of \(T_d\) is r, so \(\Xi_d\in F\). The precise positive
hexagon for the marked pair \((d,u)\), in C.30.1's convention, is
\[
                         u_0x^m u_2z^m u_1y^m=1.
 \tag{A.87.12}
\]
Its left side is \(u_0\Xi_d^{-1}u_0^{-1}\). The other required equation,
\(u(x,y)u(y,x)=1\), is already the actual canonical swap duality.
If \(\Xi_d=1\), C.30 therefore supplies the full automorphism with
\(x\mapsto x^d\), \(y\mapsto u^{-1}y^d u\), and its explicit
third-cusp conjugator. A.86 gives \(\Omega_\alpha=1\).

Conversely \(\Omega_\alpha=1\) supplies the full third-cusp equation
for these very d,u. C.29 proves that the resulting normalized
substitution is a peripheral automorphism \(\theta\). The canonical
equation gives \(\alpha\theta=\theta c\) on both generators.
By A.6 the particular peripheral conjugator \(\theta\) belongs to
\(\widehat{GT}_0\), since \(\alpha,c\) do. Its hexagon is exactly
(A.87.12), proving (A.87.3). This uses the established descent of
A.6, not merely the existence of some other gentle conjugator.

All \(u_i\) have zero abelianization and \(x+y+z=0\), so (A.87.11)
gives \(\Xi_d\in F'\) without assuming its vanishing. Also
\[
 \tau(v)=v,\qquad \tau(\delta)=v\delta^{-1}v,
 \qquad\tau(T_d)=T_d^{-1}.
\]
C.43's actual one-section theorem and \(C_F(\tau)=1\) give the unique
norm expression (A.87.4). On abelianization \(\tau\) is negative
interchange, so \((a_d)_{\rm ab}=(\ell_d,-\ell_d)\).
The norm expression vanishes exactly when \(a_d=1\); uniqueness is
not itself a vanishing argument.

\emph{3. The actual positive r-frame.}\par
The existence and uniqueness of A in (A.87.5) is the actual pure-frame
statement E.65.1, whose torsion-conjugacy and centralizer inputs are
those of E.58 and the full free-profinite-product theorem.
\citep{A_HR1985}
The exponent is +1: \(\Phi_\Sigma(r)=f^{-1}r\) has the abelianization
of r, not of \(r^{-1}\). Its conjugator can be made pure by right
multiplication by a power of r, since its S3-image centralizes r;
\(C_B(r)\cap F=1\) proves uniqueness.
Equation (A.87.5) gives \(f=\rho(A)A^{-1}\). The matrix of \(\rho\)
on \(F_{\rm ab}=\widehat{\mathbb Z}^2\) is
\[
 R=\begin{pmatrix}0&-1\\1&-1\end{pmatrix},\qquad \det(R-I)=3.
\]
Thus \((R-I)A_{\rm ab}=0\). Multiplication by3 on
\(\widehat{\mathbb Z}^2\) is injective, so the adjugate equation proves
\(A_{\rm ab}=0\), including its three-primary coordinate.

Set \(J=sfA\). Since \(Ar^{-1}A^{-1}=r^{-1}f\),
\[
 \Phi_\alpha(r)=Jr^{-1}J^{-1},\qquad
 \Phi_\alpha(J)J=1.
 \tag{A.87.13}
\]
For the second equality, the square of the first map centralizes r.
Its head \(\Phi_\alpha(J)J\) is pure, because both factors have
S3-image s, and \(C_B(r)\cap F=1\). Consequently
\(w_3=J^{-1}\alpha\) is an involution reversing r. Since
\(\alpha=v\tau\), its exact F-head is
\[
 w_3=\rho(A)^{-1}u\tau,
\]
using \(J^{-1}v=A^{-1}f^{-1}su^{-1}s=A^{-1}f^{-1}u\).
Moreover \(\Phi_\alpha(J)=s\alpha(A)\); (A.87.13) becomes
\(\iota(\alpha(A))fA=1\). Applying \(\iota\), using
\(\iota(f)=f^{-1}\), and then conjugating by \(u^{-1}\) gives
the two action identities in (A.87.6).

\emph{4. The literal cusp reflection and its normalization.}\par
The element \(\delta\alpha\) is an involution reversing \(\delta\).
Hence \(p_e=r^{-1}(\delta\alpha)r\) is an involution reversing e and z.
Keeping the semidirect multiplication order,
\[
 p_e=(r^{-1}srJr^{-1})w_3
     =z^{-1}\rho^2(A)\rho(A)^{-1}u\tau.
 \tag{A.87.14}
\]
Indeed \(r^{-1}srsr^{-1}=e^2=z^{-1}\) and
\(\rho(fA)=\rho^2(A)\), which give the second equality without
commuting any factors. The F-head in (A.87.14) has abelianization
\((1,1)\), since \(A,u\in F'\). A.86's normalization therefore
multiplies this head by z, giving
\(h_0=\rho^2(A)\rho(A)^{-1}u\) exactly. Its uniqueness follows
also directly: two zero-abelian cusp heads differ on the left by
an element of \(C_F(z)=\overline{\langle z\rangle}\), and that
procyclic group has injective abelianization in the direction
\((-1,-1)\). Finally
\[
 h_0\tau=z p_e=r^{-1}(\delta^{-1}\alpha)r,
\]
proving the rest of (A.87.6). In particular the canonical q is now
specified by the exact equation
\[
                 \rho^2(A)\rho(A)^{-1}u=q\tau(q)^{-1}.
 \tag{A.87.15}
\]
Abelianizing this equation does not determine t: the norm map for
negative interchange has the anti-diagonal as its abelian kernel.

\emph{5. Original first-cusp Alexander equation.}\par
Put \(\sigma=\operatorname{Ad}(h_0)\tau\). A.86 identifies C.44's
unique solution as
\[
 g=e(qvq^{-1}),\qquad \sigma(g)g^{-1}=z,
 \qquad g_{\rm ab}=(1-2t,2t).
\]
Conjugating both the involution and the equation by r gives
\[
 \sigma_x=\operatorname{Ad}(\delta^{-1})\Phi_\alpha|_F,
 \qquad G=rgr^{-1},\qquad \sigma_x(G)G^{-1}=x.
\]
The last equation is equivalent to (A.87.7), and uniqueness is
transported from C.44. Applying R to \(g_{\rm ab}\) gives
\(G_{\rm ab}=(-2t,1-4t)\). For an entirely original-marking formula,
write \(j_\delta=\operatorname{Ad}(\delta^{-1})|_F\). Then
\[
 \begin{gathered}
 j_\delta(x)=x,\qquad j_\delta(y)=y^{-1}x^{-1},\\
 \sigma_x(x)=x^{-1},\qquad
 \sigma_x(y)=j_\delta(f)^{-1}xyj_\delta(f).
 \end{gathered}
 \tag{A.87.16}
\]
For \(G_{\rm ab}=(b,a)\), the equation
\(\alpha(G)=j_\delta(G)x\) gives
\((-b,-a)=(b-a+1,-a)\), hence \(a=2b+1\).
Together with \(d=2k+1\) and injectivity of multiplication by2 on
\(\widehat{\mathbb Z}\), this proves (A.87.8).

\emph{6. The two actual frame coordinates and the R2 boundary.}\par
There is a unique pure \(\gamma\in F\) with
\(v=\gamma s\gamma^{-1}\): the full torsion theorem makes v
conjugate to s, its conjugator can be made pure, and
\(C_B(s)\cap F=1\) gives uniqueness. Equivalently,
\(u=\iota(\gamma)\gamma^{-1}\), so
\(\gamma_{\rm ab}=(b_0,b_0)\). In A.45's actual dihedral quotient,
\(\gamma\) therefore has even reflection parity and image
\(T_D^{\rho_D}\). The swap inverts \(T_D\), giving
\[
                k=-2\rho_D,\qquad d=1-4\rho_D.
 \tag{A.87.17}
\]
Thus \(a_{\rm reg}=d\) is equivalently \(t=\rho_D\), not an
identification of \(\rho_D\) with an ordinary abelian coordinate.
The pure s-frame \(\beta\) of \(\Phi_\alpha(s)\) also satisfies
\(\beta=\alpha(\gamma)^{-1}\gamma\), by applying
\(\Phi_\alpha\) to v and using the same centralizer. These exact
frame identities do not prove the missing coordinate equality.

The modular assignment \(s\mapsto v,\ \delta\mapsto\delta^d\)
is a homomorphism precisely when \(\Xi_d=1\). Therefore the
actual LM7\_PUSH D.50/E.58 relative lifting theorem cannot be
applied to this prospective input before its order-three relation,
or an explicitly justified corresponding relative relation, is
established. In the pants formulation the same missing input is
the specified third-peripheral conjugator for the given first two
images. C.44's existing relative Alexander automorphism is not
asserted to preserve those three markings. The results above
retain every actual frame and reduce the coefficient to (A.87.8),
but supply neither that equality nor the nonlinear vanishing in
(A.87.3). This completes the claimed scoped comparison.
\end{proof}
